\subsection{Divisores y normalización}
\label{subsection:IV.21.8}

\begin{lemma}[21.8.1]
\label{IV.21.8.1}
Sea $f:X'\to X$ un morfismo entero de preesquemas.
Para todo $\sh{O}_{X'}$-Módulo localmente libre $E'$ de rango constante $n$ y todo $x\in X$, existe un entorno abierto $U$ de $x$ en $X$ tal que, si se pone $U'=f^{-1}(U)$, $E'|U'$ sea isomorfo a $\sh{O}_{U'}^n$.
\end{lemma}

\begin{proof}
Como la cuestión es local sobre $X$, podemos limitarnos al caso en que $X=\Spec(A)$ es afín, en cuyo caso se tiene $X'=\Spec(A')$, donde $A'$ es una $A$-álgebra entera.
Entonces $A'$ es límite inductivo de sus sub-$A$-álgebras \emph{finitas} $A'_\lambda$.
Poniendo $X'_\lambda=\Spec(A'_\lambda)$ y denotando por $p_\lambda:X'\to X'_\lambda$ el morfismo estructural, se deduce de \sref{IV.8.5.2} y \sref{IV.8.5.5} que hay un índice $\lambda$ y un $\sh{O}_{X'_\lambda}$-Módulo localmente libre $E'_\lambda$ de rango $n$ tal que $E'=p_\lambda^*(E'_\lambda)$, y bastará evidentemente demostrar el lema para $X'_\lambda$ y $E'_\lambda$;
en otras palabras, podemos limitarnos al caso en que el morfismo $f$ es \emph{finito}.
Pongamos $B=\sh{O}_{X,x}$ y sea $B'$ el anillo del esquema afín $X'_0=X'\times_X\Spec(\sh{O}_{X,x})$;
como $B$ es un anillo local y $B'$ es una $B$-álgebra finita, $B'$ es un anillo semilocal (Bourbaki, \emph{Alg.\ comm.}, cap.~V, \textsection~2, n\textsuperscript{o}~1, prop.~3);
se concluye que el $\sh{O}_{X'_0}$-Módulo localmente libre $E'_0=E'\otimes_{\sh{O}_{X'}}\sh{O}_{X'_0}$ es isomorfo a $\sh{O}_{X'_0}^n$ (Bourbaki, \emph{Alg.\ comm.}, cap.~II, \textsection~5, n\textsuperscript{o}~3, prop.~5).
Considerando $X'_0$ como límite
\oldpage[IV-4]{281}
proyectivo de los preesquemas inducidos por $X'$ sobre los abiertos $f^{-1}(U)$, donde $U$ recorre el conjunto filtrante de los entornos abiertos afines de $x$ en $X$, siguiendo el método de \sref{IV.8.1.2}, \emph{a)}), y aplicando \sref{IV.8.5.2.5}, se obtiene la conclusión buscada.
\end{proof}

\begin{corollary}[21.8.2]
\label{IV.21.8.2}
Sea $f:X'\to X$ un morfismo entero de preesquemas.
Entonces:
\begin{enumerate}[label=(\roman*)]
  \item Se tiene $\RR^1 f_*(\sh{O}^*_{X'})=0$ \sref[0]{0.12.2.1}.
  \item El grupo $\operatorname{Pic}(X')$ es canónicamente isomorfo a $\HH^1(X,f_*(\sh{O}^*_{X'}))$.
\end{enumerate}
\end{corollary}

\begin{proof}
\begin{enumerate}[label=(\roman*)]
  \item $\RR^1 f_*(\sh{O}^*_{X'})$ es el haz asociado al prehaz de grupos conmutativos $U\mapsto\HH^1(f^{-1}(U),\sh{O}^*_{X'})$ sobre $X$ (\emph{loc.\ cit.});
la fibra de este haz en un punto $x$ puede, pues, identificarse \sref[0]{0.5.4.7} con el grupo conmutativo $\varinjlim\operatorname{Pic}(f^{-1}(U))$, donde $U$ recorre el conjunto de los entornos abiertos de $x$, y los homomorfismos de transición $\operatorname{Pic}(f^{-1}(U'))\to\operatorname{Pic}(f^{-1}(U))$ para dos abiertos $U\subset U'$ están definidos por \sref{IV.21.3.2.4}.
Ahora bien, para todo $x\in X$, todo entorno abierto $U'$ de $x$ en $X$ y todo elemento $\zeta\in\operatorname{Pic}(f^{-1}(U'))$, se deduce de \sref{IV.21.8.1} que existe un entorno abierto $U\subset U'$ de $x$ tal que la imagen de $\zeta$ en $\operatorname{Pic}(f^{-1}(U))$ sea nula.
Por tanto, la fibra de $\RR^1f_*(\sh{O}^*_{X'})$ en el punto $x$ es nula.

  \item La sucesión espectral de Leray para el functor compuesto $\sh{F}'\mapsto\Gamma(X,f_*(\sh{F}'))$ de haces de grupos conmutativos sobre $X'$ (T, 2.4) da la sucesión exacta de términos de grado bajo
\[
  0\to\HH^1(X,f_*(\sh{O}^*_{X'}))\to\HH^1(X',\sh{O}^*_{X'})\to\HH^0(X,\RR^1f_*(\sh{O}^*_{X'}))
\]
y la conclusión se deduce de (i) y del isomorfismo entre $\operatorname{Pic}(X')$ y $\HH^1(X',\sh{O}^*_{X'})$ \sref[0]{0.5.4.7}.
\end{enumerate}
\end{proof}

\begin{proposition}[21.8.3]
\label{IV.21.8.3}
Sea $f=(\psi,\theta):X'\to X$ un morfismo entero de preesquemas;
supongamos que, para todo abierto $U$ de $X$, el homomorfismo $\Gamma(\theta):\Gamma(U,\sh{O}_X)\to\Gamma(U,f_*(\sh{O}_{X'}))$ transforma los elementos regulares en elementos regulares, lo que permite prolongar canónicamente el homomorfismo $\theta:\sh{O}_X\to f_*(\sh{O}_{X'})$ a un homomorfismo de haces de anillos $\theta':\sh{M}_X\to f_*(\sh{M}_{X'})$, de donde se obtienen homomorfismos de haces de grupos multiplicativos $\theta^*:\sh{O}^*_X\to f_*(\sh{O}^*_{X'})$ y $\theta'^*:\sh{M}^*_X\to f_*(\sh{M}^*_{X'})$, de los que se obtiene por paso a los cocientes un homomorfismo $\theta''^*:\sh{D}iv_X\to f_*(\sh{D}iv_{X'})$.
Se tiene entonces un diagrama conmutativo
\[
\label{IV.21.8.3.1}
  \xymatrix{
    1 \ar[r] & \sh{O}^*_X \ar[r] \ar[d]^{\theta^*} & \sh{M}^*_X \ar[r] \ar[d]^{\theta'^*} & \sh{D}iv_X \ar[r] \ar[d]^{\theta''^*} & 0\\
    1 \ar[r] & f_*(\sh{O}^*_{X'}) \ar[r] & f_*(\sh{M}^*_{X'}) \ar[r] & f_*(\sh{D}iv_{X'}) \ar[r] & 0
  }
  \tag{21.8.3.1}
\]
en el que las dos filas son exactas.
\end{proposition}

\begin{proof}
El único punto que hay que verificar es la exactitud de la segunda fila del diagrama, que se deduce de aplicar la sucesión exacta de cohomología para el functor $f_*$ a la sucesión exacta de haces de grupos conmutativos sobre $X'$
\[
  1\to\sh{O}^*_{X'}\to\sh{M}^*_{X'}\to\sh{D}iv_{X'}\to0
\]
y de \sref{IV.21.8.2}.
\end{proof}

\begin{corollary}[21.8.4]
\label{IV.21.8.4}
Si, además de las hipótesis de \sref{IV.21.8.3}, el homomorfismo $\theta'$ es un isomorfismo de haces de anillos, entonces $\theta^*:\sh{O}^*_X\to f_*(\sh{O}^*_{X'})$ es inyectivo, $\theta''^*:\sh{D}iv_X\to f_*(\sh{D}iv_{X'})$ es sobreyectivo y $\operatorname{Ker}(\theta''^*)$ es isomorfo a $\operatorname{Coker}(\theta^*)$.
\end{corollary}

\begin{proof}
Es una consecuencia inmediata del diagrama de la serpiente (Bourbaki, \emph{Alg.\ comm.}, cap.~I, \textsection~2, n\textsuperscript{o}~4, prop.~2) aplicado al diagrama \sref{IV.21.8.3.1}.
\end{proof}

\oldpage[IV-4]{282}

\begin{proposition}[21.8.5]
\label{IV.21.8.5}
Sean $X$, $X'$ dos preesquemas localmente noetherianos y $f:X'\to X$ un morfismo entero;
se supone que existe un abierto esquemáticamente denso $U$ en $X$ tal que $f^{-1}(U)$ sea esquemáticamente denso en $X'$ (cf.\ \sref{IV.20.3.5}), y que el morfismo $f^{-1}(U)\to U$, restricción de $f$, sea un isomorfismo.
Entonces:
\begin{enumerate}[label=(\roman*)]
  \item El homomorfismo $\theta':\sh{M}_X\to f_*(\sh{M}_{X'})$ está definido y es biyectivo, el homomorfismo
\[
  \theta^*:\sh{O}^*_X\to f_*(\sh{O}^*_{X'})
\]
es inyectivo, el homomorfismo $\theta''^*:\sh{D}iv_X\to f_*(\sh{D}iv_{X'})$ es sobreyectivo, $\operatorname{Ker}(\theta''^*)$ es isomorfo a $\sh{N}=\operatorname{Coker}(\theta^*)$ y el soporte del haz de grupos multiplicativos $\sh{N}$ está contenido en $S=X-U$.
  \item Supongamos además que el conjunto $S$ es discreto y pongamos, para abreviar, $\sh{O}'_X=f_*(\sh{O}_{X'})$.
Entonces se tiene un diagrama conmutativo
\[
\label{IV.21.8.5.1}
  \xymatrix{
    1 \ar[r] & \prod_{s\in S}\sh{O}'^{*}_{X,s}/\sh{O}^*_{X,s} \ar[r]^-{j} \ar[d] & \operatorname{Div}(X) \ar[r] \ar[d]^{k_X} & \operatorname{Div}(X') \ar[r] \ar[d]^{k_{X'}} & 0\\
    1 \ar[r] & \bigl(\prod_{s\in S}\sh{O}'^{*}_{X,s}/\sh{O}^*_{X,s}\bigr)/\operatorname{Im}\Gamma(X',\sh{O}^*_{X'}) \ar[r]^-{j'} & \operatorname{Pic}(X) \ar[r] & \operatorname{Pic}(X') \ar[r] & 0
  }
  \tag{21.8.5.1}
\]
donde las filas son exactas y la flecha vertical de la izquierda es el homomorfismo canónico.
\end{enumerate}
\end{proposition}

\begin{proof}
\begin{enumerate}[label=(\roman*)]
  \item La hipótesis implica que se tiene un isomorfismo canónico $\sh{M}''_X\simto f_*(\sh{M}''_{X'})$ para los haces de gérmenes de pseudo-funciones \sref{IV.20.2.2}, de donde se deduce la afirmación relativa a $\theta'$, dada la existencia de los isomorfismos canónicos $\sh{M}_X\simto\sh{M}''_X$ y $\sh{M}_{X'}\simto\sh{M}''_{X'}$ \sref{IV.20.2.11}.
El resto de la afirmación (i) se deduce de \sref{IV.21.8.4}, salvo lo relativo al soporte de $\sh{N}$, que resulta directamente de la hipótesis sobre $U$.

  \item Aplicando la sucesión exacta de cohomología a las dos sucesiones exactas de haces de grupos conmutativos
\[
\label{IV.21.8.5.2}
\begin{gathered}
  1\to\sh{N}\to\sh{D}iv_X\to f_*(\sh{D}iv_{X'})\to0,\\
  1\to\sh{O}^*_X\to f_*(\sh{O}^*_{X'})\to\sh{N}\to1,
\end{gathered}
  \tag{21.8.5.2}
\]
se obtienen respectivamente las sucesiones exactas
\[
\label{IV.21.8.5.3}
\begin{gathered}
  1\to\Gamma(X,\sh{N})\xrightarrow{j}\operatorname{Div}(X)\to\operatorname{Div}(X')\to\HH^1(X,\sh{N}),\\
  \Gamma(X',\sh{O}^*_{X'})\to\Gamma(X,\sh{N})\xrightarrow{\partial}\operatorname{Pic}(X)\to\operatorname{Pic}(X')\to\HH^1(X,\sh{N}),
\end{gathered}
  \tag{21.8.5.3}
\]
\oldpage[IV-4]{283}
teniendo en cuenta el isomorfismo canónico $\operatorname{Pic}(X')\simto\HH^1(X,f_*(\sh{O}^*_{X'}))$ \sref{IV.21.8.2}.
Ahora bien, como el soporte de $\sh{N}$ está contenido en el conjunto discreto $S$, cerrado en $X$, se tiene
\[
  \HH^1(X,\sh{N})=\HH^1(S,\sh{N}|S)\quad\text{(G, II, 4.9.2)},\qquad
  \HH^1(S,\sh{N}|S)=\prod_{s\in S}\HH^1(\{s\},\sh{N}_s)=0
\]
por definición de la cohomología (G, II, 4.4).
Asimismo, se tiene $\Gamma(X,\sh{N})=\prod_{s\in S}\sh{N}_s$ y $\sh{N}_s=\sh{O}'^{*}_{X,s}/\sh{O}^*_{X,s}$ en virtud de la segunda sucesión exacta \sref{IV.21.8.5.2}.
Queda precisar las inyecciones $j$ y $j'$.
Ahora bien, se puede describir una sección $t$ de $\sh{N}$ sobre $X$ tomando un recubrimiento de $X$ formado por $U$ y por abiertos $U_i$ tales que $U_i\cap S$ se reduzca a un solo punto $s_i$, y que $U_i\cap U_j\subset U$ para $i\neq j$, y considerando después, para cada $i$, una sección $t_i$ de $\sh{N}$ sobre $U_i$, necesariamente tal que $(t_i)_{s_i}\in\sh{O}'^{*}_{X,s_i}/\sh{O}^*_{X,s_i}$ y $(t_i)_x=0$ en los puntos $x\in U_i-\{s_i\}$.
El germen $(t_i)_{s_i}$ proviene de una sección $u_i$ de $\sh{O}^*_{X'}$ sobre $f^{-1}(U_i)$, que también puede considerarse como una sección de $\sh{M}^*_{X'}$ sobre $f^{-1}(U_i)$ y, por tanto, como una sección de $\sh{M}^*_X$ sobre $U_i$ (en virtud de la hipótesis);
a esta sección le corresponde canónicamente una sección $d_i\in\Gamma(U_i,\sh{D}iv_X)$ y, como en $f^{-1}(U\cap U_i)$ la restricción de $u_i$ se identifica con una sección de $\sh{O}^*_X$ sobre $U\cap U_i$ en virtud de la hipótesis, las restricciones de los $d_i$ a $U\cap U_i$ son todas nulas; por tanto, los $d_i$ son las restricciones de un mismo divisor perteneciente a $\operatorname{Div}(X)$, que es la imagen de la sección $t$ por $j$.
Asimismo, la imagen de $t$ por $\partial:\Gamma(X,\sh{N})\to\HH^1(X,\sh{O}^*_X)$ proviene del cociclo igual a la restricción de $u_i$ en $U\cap U_i$ para cada $i$, y a $(u_i|(U_i\cap U_j))(u_j|(U_i\cap U_j))^{-1}$ en $U_i\cap U_j$, de donde se obtienen la expresión de $j'(t)$ y la conmutatividad del diagrama \sref{IV.21.8.5.1}.
\end{enumerate}
\end{proof}

\begin{remark}[21.8.6]
\label{IV.21.8.6}
\begin{enumerate}[label=(\roman*)]
  \item Las condiciones de aplicación de \sref{IV.21.8.5}, (i) se satisfacen en particular cuando $X$ es un preesquema localmente noetheriano reducido que solo tiene un número finito de componentes irreducibles, $X'$ es su normalización y $X'$ es también localmente noetheriano (\sref[II]{II.6.3.8});
el $\sh{O}_X$-Módulo $\sh{D}iv_X$ es entonces una \emph{extensión} de $f_*(\sh{D}iv_{X'})$ por el conúcleo $\sh{N}$ de $\sh{O}^{*}_X\to f_*(\sh{O}^*_{X'})$, que puede considerarse conocido.
Si además $X$ es una curva algebraica (reducida) sobre un cuerpo $k$, se cumplen las condiciones de aplicación de \sref{IV.21.8.5}, (ii).

  \item Cuando se cumplen las condiciones de aplicación de \sref{IV.21.8.5}, (ii), y además el homomorfismo canónico global $\Gamma(X,\sh{O}_X)\to\Gamma(X',\sh{O}_{X'})$ es \emph{biyectivo}, se ve que los núcleos de los homomorfismos sobreyectivos $\operatorname{Div}(X)\to\operatorname{Div}(X')$ y $\operatorname{Pic}(X)\to\operatorname{Pic}(X')$ son \emph{isomorfos}.

  \item Cuando se cumplen las condiciones de aplicación de \sref{IV.21.8.5}, (ii), se ve que el homomorfismo $\operatorname{Div}(X)\to\operatorname{Div}(X')$ solo puede ser inyectivo (en cuyo caso es biyectivo) si $\sh{O}'^{*}_{X,s}=\sh{O}^*_{X,s}$ para todo $s\in S$.
Para un $s\in S$ tal que el anillo $\sh{O}'_{X,s}$ sea \emph{local} (lo que, teniendo en cuenta el hecho de que $X'=\Spec(\sh{O}'_X)$ \sref[II]{II.1.3}, significa que solo existe un punto $s'\in X'$ sobre $s$), esto implica necesariamente que los cuerpos residuales $\boldsymbol{k}(s)$ y $\boldsymbol{k}(s')$ son iguales y que $1+\mathfrak{m}_s=1+\mathfrak{m}_{s'}$, y por tanto equivale finalmente a la relación $\sh{O}'_{X,s}=\sh{O}_{X,s}$, o también (teniendo en cuenta la hipótesis) al hecho de que existe un entorno abierto $V$ de $s$ en $X$ tal que el morfismo $f^{-1}(V)\to V$ sea un isomorfismo.
En el caso general, el anillo $\sh{O}'_{X,s}$ es semilocal (esto es evidente cuando el morfismo $f$ es finito y, en el caso general, se deduce de \sref[0]{0.23.2.6}); el homomorfismo canónico $\sh{O}^{*}_{X,s}\to\sh{O}'^{*}_{X,s}$ define, por paso a los cocientes, un homomorfismo del grupo multiplicativo $\boldsymbol{k}(s)^*$ en el producto de los grupos multiplicativos $\boldsymbol{k}(s'_j)^*$, siendo los $s'_j$ los puntos (en
\oldpage[IV-4]{284}
número $>1$) de $X'$ situados sobre $s$.
Es inmediato que este homomorfismo solo puede ser biyectivo si $\boldsymbol{k}(s)$ y todos los $\boldsymbol{k}(s'_j)$ son iguales al cuerpo $\bb{F}_2$ de $2$ elementos;
además, si se cumple esta condición, es necesario también que el grupo multiplicativo $1+\mathfrak{m}_s$ tenga por imagen $1+\mathfrak{r}$, donde $\mathfrak{r}$ es el radical de $\sh{O}'_{X,s}$, o, equivalentemente, que $\mathfrak{m}_s=\mathfrak{r}$, lo que implica que $\sh{O}'_{X,s}\otimes_{\sh{O}_{X,s}}\boldsymbol{k}(s)$ es un producto directo de cuerpos isomorfos a $\boldsymbol{k}(s)$ (lo que, cuando el morfismo $f$ es finito, implica que no está ramificado en los puntos $s'_j$ \sref{IV.17.4.1}).
Si, por ejemplo, \emph{ningún} cuerpo residual de $X$ es isomorfo a $\bb{F}_2$, el homomorfismo canónico $\operatorname{Div}(X)\to\operatorname{Div}(X')$ solo puede ser biyectivo si $f$ es un \emph{isomorfismo}.
En el caso en que el homomorfismo canónico $\Gamma(X,\sh{O}_X)\to\Gamma(X',\sh{O}_{X'})$ es biyectivo, se concluye de lo anterior y de (ii) que, en las consideraciones precedentes, se puede reemplazar el homomorfismo $\operatorname{Div}(X)\to\operatorname{Div}(X')$ por el homomorfismo $\operatorname{Pic}(X)\to\operatorname{Pic}(X')$.
\end{enumerate}
\end{remark}
